\section{Null layers with a relative frame}
\label{sec:null}

The next obstruction will resolve the null-vector ties in the harmonic
filtration.  The frame $f$ need not neighbor every operated color, so its
fixed generator alone does not provide a boundary coordinate for a given
color set.  We instead use the connected support of an individual family
member.  This gives the required obstruction at every integer level.

\begin{proposition}[Relative null-layer obstruction]
\label{null:obstruction}
Let $I\subseteq S$ be nonempty and connected, let $c_I>0$ satisfy
$C_Ic_I=0$, and extend $c$ by zero outside $I$.  Let $J\subseteq I$ be
spherical, possibly empty, let $q\in\mathbb Z$, and fix $A_Jx$.
Suppose its fixed data on $(I\setminus J)\cup\partial I$ are capped at
$q$ and satisfy
\begin{equation}
 m^q\big|_{(I\setminus J)\cup\partial I}
   =c\big|_{(I\setminus J)\cup\partial I}.
 \label{null:exterior}
\end{equation}
It is impossible that every representative $z\in A_Jx$ capped at $q$
on $I$ satisfies
\begin{equation}
 m_j^q(z)>c_j\qquad\text{for every }j\in J.
 \label{null:strict}
\end{equation}
Thus, if $J\ne\varnothing$, some such representative has
$m_j^q(z)\le c_j$ for at least one $j\in J$.  If $J=\varnothing$,
the exterior hypotheses themselves are impossible.
\end{proposition}

We first develop the linear Hom calculation used in the proof.  All
complexes, maps, and pairings retain both gradings, and all differential
signs use cochain degree.

\subsection{Linear Hom and its permanent middle column}

For pure projective layers $V,W$ of levels $l,k$, respectively, let
$\mathcal L_j^{n,v}(V,W)$ be the space of maps of path length $j$,
cochain degree $n$, and internal degree $v$.  It is zero unless
$j\in\{0,1,2\}$ and
\begin{equation}
 n+v=k-l+j.
 \label{null:hom-degree}
\end{equation}
Indeed, between generator placements $(a,u)$ and $(a',u')$, these
degrees are $n=a'-a$ and $v=u'+j-u$.  Write $\delta_V,\delta_W$ for
the linear differentials.  The differential
\[
 D_1g=\delta_Wg-(-1)^n g\delta_V
\]
raises path length and cochain degree by one and preserves internal
degree.  Denote its cohomology in the displayed space by
$\mathsf H_j^{n,v}(V,W)$, and put
\[
 h_j(V,W)=\sum_{n,v}\dim_{\kk}\mathsf H_j^{n,v}(V,W).
\]
The sums are finite.  In length zero there are no incoming boundaries:
a cocycle is a homogeneous map of coefficient arrays commuting with
their arrows up to its cochain sign.  Its image is a graded
subrepresentation, with the corresponding shift in placements.

\begin{lemma}[Pairing and middle vanishing]
\label{null:hom}
For pure layers, including layers obtained from graded
subrepresentations, there are perfect pairings
\begin{equation}
 \mathsf H_j^{n,v}(V,W)\ \times\
 \mathsf H_{2-j}^{-n,\,2-v}(W,V)\longrightarrow\kk.
 \label{null:pairing}
\end{equation}
In particular,
\begin{equation}
 \begin{split}
 h_2(V,W)&=h_0(W,V),\\
 (\dim V)^{\mathsf T}\widetilde C\,\dim W
   &=h_0(V,W)+h_0(W,V)-h_1(V,W),
 \end{split}
 \label{null:euler}
\end{equation}
where dimensions count projective generators, summed over placements.
For every $z\in A$, every pair of labels $L,H$, and all levels $l,k$,
\begin{equation}
 h_1\bigl(Y^l(E_L(z)),Y^k(E_H(z))\bigr)=0.
 \label{null:middle}
\end{equation}
The vanishing persists after replacing either layer by a direct
summand as a pure layer.
\end{lemma}

\begin{proof}
The symmetric Frobenius trace pairs opposite paths of complementary
lengths, with total internal degree two.  On a complex use the
cochain supertrace
\[
 \operatorname{Str}_W(a)
   =\sum_c(-1)^c\tau\bigl(\operatorname{tr}_{\mathrm{mat}}(a|_{W^c})\bigr),
\]
where the matrix trace sums diagonal path entries and $\tau$ is the
Frobenius trace, not the ordinary trace of the underlying linear map.  Composition
and this trace give perfect pairings between the corresponding length
spaces in Equation~\eqref{null:pairing}.  Cyclicity gives
\[
 \langle D_1g,h\rangle+(-1)^n\langle g,D_1h\rangle=0
\]
in complementary degrees.  Finite-dimensional duality yields the
pairings on cohomology.  Units and tokens each contribute the
equal-vertex pairing to the length Euler characteristic; arrows
contribute minus the adjacency multiplicities.  This proves
Equation~\eqref{null:euler}.  The argument only uses the linear
differentials and therefore also applies to subrepresentations.

For Equation~\eqref{null:middle}, take minimal models of $E_L(z)$ and
$E_H(z)$ and filter their full bigraded Hom complex by paths of length
at least $j$.  Its differential is $D_1+D_2$, with $D_a$ raising path
length by $a$.  The first nonzero differential is $D_1$, whose
cohomology splits by the ordered source and target layer levels.
There are only three columns.  After this linear page, the only
possible differential goes from length zero to length two; every
middle-column class is therefore permanent.

The action in Proposition~\ref{cat:action} preserves actual bigraded
Hom cohomology.  Between the initial projectives $P_{f,L},P_{f,H}$
it is zero for $L\ne H$; for $L=H$ it is one-dimensional in each of
bidegrees $(0,0)$ and $(0,2)$ and zero elsewhere.  Both diagonal
dimensions already have permanent representatives.  The identity of
the full complex is a nonzero length-zero class: a boundary has no
length-zero component.  For the other class, put a normalized token
on one diagonal projective summand of a nonzero minimal model and put
zero on the other entries.  This is a cycle of bidegree $(0,2)$,
because a token multiplied by a radical entry is zero.  Its
supertrace is $(-1)^c\ne0$ if that summand lies in cochain placement
$c$, whereas every Hom boundary has zero supertrace.  It is thus a
nonzero actual class in filtration two; filtration three is zero.

The identity and token exhaust the actual diagonal dimensions, and
the off-diagonal cohomology vanishes.  Hence every permanent middle
class is zero, proving Equation~\eqref{null:middle}.  A direct
summand of a pure layer gives a direct summand of its linear Hom
complex, so the same vanishing holds for it.  This last assertion
does not require the summand to split off the full transported complex.
\end{proof}

For finite label-indexed families $\mathcal V=(V_L)_L$ and
$\mathcal W=(W_H)_H$, define the weighted dimensions
\begin{equation}
 h_j^{\mathrm{wt}}(\mathcal V,\mathcal W)
    =\sum_{L,H}w_Lw_H h_j(V_L,W_H).
 \label{null:weighted-dimension}
\end{equation}
They need not be integers.  If the two weighted generator arrays are
$\phi v$ and $\phi y$, then
\[
 \sum_{L,H}w_Lw_H(\dim V_L)^{\mathsf T}\widetilde C\dim W_H
   =D_0v^{\mathsf T}Cy.
\]
Equation~\eqref{null:euler} therefore becomes
\begin{equation}
 D_0v^{\mathsf T}Cy
  =h_0^{\mathrm{wt}}(\mathcal V,\mathcal W)
   +h_0^{\mathrm{wt}}(\mathcal W,\mathcal V)
   -h_1^{\mathrm{wt}}(\mathcal V,\mathcal W).
 \label{null:weighted-euler}
\end{equation}

\subsection{The capped secondary injection}

Fix a representative $z$, and abbreviate
$\mathcal Y^d=(Y^d(E_L(z)))_L$.  The following lemma uses a cap to
remove one of the two possible blocks of the secondary Hom differential.

\begin{lemma}[Secondary injection]
\label{null:injection}
Suppose the family at $z$ is capped at $q$ on $I\subseteq S$.
For each $L$, let $M_L$ be a direct summand of $Y^q(E_L(z))$ as a
pure layer, supported on $\widetilde I$, and put
$\mathcal M=(M_L)_L$.  For $d<q$ there are injections
\begin{equation}
 \mathsf H_0^{n,v}\bigl(M_L,Y^d(E_H(z))\bigr)
 \lhook\joinrel\longrightarrow
 \mathsf H_2^{n+1,v}\bigl(M_L,Y^{d-1}(E_H(z))\bigr),
 \qquad n+v=d-q.
 \label{null:secondary-map}
\end{equation}
Consequently,
\begin{equation}
 h_0^{\mathrm{wt}}(\mathcal M,\mathcal Y^d)
    \le h_0^{\mathrm{wt}}(\mathcal Y^{d-1},\mathcal M).
 \label{null:secondary-bound}
\end{equation}
\end{lemma}

\begin{proof}
Extend a length-zero linear cocycle $g:M_L\to Y^d(E_H(z))$ by zero
on the other source summands.  Because $M_L$ is a direct summand for
the linear differential, this embeds its class into the linear page
of $\Hom(E_L(z),E_H(z))$.  No length-zero boundary is present.
The secondary differential is represented by
\[
 d_{E_H}^{(2)}g-(-1)^n g d_{E_L}^{(2)},
\]
where the superscript denotes the length-two part of the differential.
The two possible ordered layer pairs of its image are
\[
 (q,d-1)\quad\text{and}\quad(q+1,d),
\]
respectively.  The latter term would require a level-$(q+1)$ source
at a vertex of $\widetilde I$, since a length-two entry is a token
at a single vertex.  The cap excludes it.  Thus the secondary image
lies in the first block and gives Equation~\eqref{null:secondary-map}.

If this image were zero, the length-zero class would be permanent:
there are no incoming differentials and no later possible outgoing
differential.  Its total degree is $n+v=d-q<0$, but actual Hom
cohomology between the transported family members vanishes in every
negative total degree by the initial projective calculation in
Lemma~\ref{null:hom}.  The class is therefore zero.  Finally, sum
the resulting dimension inequalities with weights $w_Lw_H$ and use
Equation~\eqref{null:pairing}, interchanging the two label indices.
This proves Equation~\eqref{null:secondary-bound}.
\end{proof}

\subsection{Convex separation from a null diagram}

The convex-hull conclusion of Proposition~\ref{cap:hull} interacts
with the Hom calculation through the sign of a null vector at its
exterior neighbors.

\begin{lemma}[Separation]
\label{null:separation}
Let $I\subseteq S$ be nonempty and connected, with $c_I>0$ and
$C_Ic_I=0$, and extend $c$ by zero.  Let $J\subseteq I$.
Suppose $\mathcal M=(M_L)_L$ consists of direct summands of a single
layer family $\mathcal Y^q$, supported on $\widetilde I$, with
weighted generator array $\phi m$.  Assume $m-c$ is supported on
$J$, and $(m-c)_J$ is an interior point in $\mathbb R^J$ of the
convex hull of finitely many vectors $s_J$.  Each such $s$ is zero
off $J$, and $\phi s$ is the weighted generator array of a family
of graded subrepresentations $\mathcal S\subseteq\mathcal M$
supported on $\widetilde J$.

If the socle of $\mathcal Y^d$ vanishes on $J$, and its folded
weighted vector is $y=m^d(z)$, then
\begin{equation}
 \begin{gathered}
 y|_{\partial I}=0,\qquad (Cy)_J=0,\\
 h_0^{\mathrm{wt}}(\mathcal M,\mathcal Y^d)
   =h_0^{\mathrm{wt}}(\mathcal Y^d,\mathcal M)=0.
 \end{gathered}
 \label{null:separation-conclusions}
\end{equation}
The usual zero-dimensional interior convention is used when
$J=\varnothing$.
\end{lemma}

\begin{proof}
Write $\mathcal Y=\mathcal Y^d$, and fix one of the subfamilies
$\mathcal S$.  A nonzero homogeneous length-zero cocycle from a
member of $\mathcal S$ to one of $\mathcal Y$ has a nonzero graded
image supported on $\widetilde J$.  That image has a socle: in its
highest occupied cochain placement every outgoing coefficient arrow
is zero, since arrows raise that placement.  This gives a socle of
the target on $J$, contrary to the hypothesis.  Hence
\[
 h_0^{\mathrm{wt}}(\mathcal S,\mathcal Y)=0.
\]
Inclusion of $\mathcal S$ in $\mathcal M$ also gives
\[
 h_0^{\mathrm{wt}}(\mathcal Y,\mathcal S)
   \le h_0^{\mathrm{wt}}(\mathcal Y,\mathcal M),
\]
because length-zero cocycles have no boundaries to quotient out.
Apply Equation~\eqref{null:weighted-euler} to these two families.
Only nonnegativity of $h_1^{\mathrm{wt}}(\mathcal S,\mathcal Y)$
is needed; for $\mathcal M$, middle vanishing does hold by
Lemma~\ref{null:hom}.  Subtraction gives
\begin{equation}
 D_0(m-s)^{\mathsf T}Cy
     \ge h_0^{\mathrm{wt}}(\mathcal M,\mathcal Y)\ge0.
 \label{null:convex-inequality}
\end{equation}

This inequality is affine-linear in $s$, so it holds on the entire
hull.  At its interior point $e=m-c$, its left side is
$D_0c^{\mathsf T}Cy$.  Symmetry of $C$ and the null equation give
\begin{equation}
 c^{\mathsf T}Cy
    =-\sum_{a\notin I}\sum_{i\in I}b_{ai}c_i y_a\le0.
 \label{null:boundary-sign}
\end{equation}
It is strictly negative if $y$ meets $\partial I$.  Comparing with
Equation~\eqref{null:convex-inequality} proves boundary vanishing,
$c^{\mathsf T}Cy=0$, and
$h_0^{\mathrm{wt}}(\mathcal M,\mathcal Y)=0$.

For each $j\in J$, ordinary interior membership supplies the two
nearby hull points $e\pm\varepsilon\mathbf e_j$.  Their inequalities,
whose value at $e$ is zero, force $(Cy)_j=0$.  Therefore
\[
 m^{\mathsf T}Cy
   =c^{\mathsf T}Cy+(m-c)_J^{\mathsf T}(Cy)_J=0.
\]
The Euler identity for $(\mathcal M,\mathcal Y)$ now gives the
remaining Hom vanishing.  When $J$ is empty, $m=c$ and no coordinate
variation is required.
\end{proof}

\subsection{Proof of the obstruction}

\begin{proof}[Proof of Proposition~\ref{null:obstruction}]
Assume Equation~\eqref{null:strict} holds at every representative
capped on $I$.  Since
$\partial J\subseteq(I\setminus J)\cup\partial I$, cap accessibility
in Lemma~\ref{cap:access} supplies a representative capped on $J$
and hence on all of $I$.  We will use only cap-preserving moves in
$J$, so the exterior data in Equation~\eqref{null:exterior} remain
fixed throughout.

\emph{The convex data at each capped representative.}
Let $M_L$ be the part of $Y^q(E_L(z))$ supported on $\widetilde I$.
It is a direct summand of the pure layer, because the level-$q$
coefficients on $\partial I$ vanish.  This does not assert a splitting
of $E_L(z)$.  Its weighted generator array is $\phi m$, where $m$
agrees with $m^q(z)$ on $I$ and is zero elsewhere.  In particular,
$m-c$ is supported on $J$.

The harmonic prescription on $J$ is $c_J$: its boundary data are
those of $c$, and $C_Ic_I=0$.  Remote coordinates outside
$I\cup\partial I$ do not enter those equations.  The strict part of
Proposition~\ref{cap:hull} and the universal hypothesis
Equation~\eqref{null:strict} therefore give
\begin{equation}
 (m-c)_J\in\operatorname{int}_{\mathbb R^J}
       \operatorname{conv}\{s_J\},
 \label{null:interior}
\end{equation}
with the subrepresentation families required in Lemma~\ref{null:separation}.
They lie in $\mathcal M$ because their support is in $\widetilde J$.
The assertion is renewed after every capped move: both $\mathcal M$
and its hull may change, but the strict hypothesis is quantified over
every capped representative.

\emph{An occupied boundary below $q$.}
The set $I\setminus J$ is nonempty, since $C_J$ is positive definite
whereas $C_I$ has a positive null vector.  Choose a color there.
Its prescribed positive level-$q$ multiplicity occurs in some member
$E_L(z)$.  By the family properties in Proposition~\ref{cat:family},
the minimal color support of that member is connected and contains
$f$, which lies outside $I$.  A support path from the chosen color
to $f$ meets $\partial I$.  There is consequently an occupied
boundary level.  Let $d_b$ be the lowest such level over the family.
The boundary cap and zero boundary data at level $q$ give
\begin{equation}
 d_b<q.
 \label{null:boundary-level}
\end{equation}
All boundary colors lie outside $J$, so $d_b$ is fixed in the coset.
No adjacency of $f$ to $I$, and no sign restriction on $q$, was used.

\emph{The upward induction.}
Choose $d_{\min}$ below or equal to every initially occupied layer.
Positive amplitude preserves this lower bound.  At the beginning of
a stage $d\in\{d_{\min},\ldots,d_b\}$, maintain, for every $k<d$,
\begin{equation}
 \begin{gathered}
 \soc_{\widetilde J}\mathcal Y^k=0,\qquad (Cm^k)_J=0,\\
 h_0^{\mathrm{wt}}(\mathcal M,\mathcal Y^k)
   =h_0^{\mathrm{wt}}(\mathcal Y^k,\mathcal M)=0.
 \end{gathered}
 \label{null:induction}
\end{equation}
Here $\mathcal M$ always refers to the current level-$q$ summand.
Initially these conditions are vacuous below the lower bound.
Since $d\le d_b<q$, Lemma~\ref{null:injection} and the condition at
$d-1$ imply
\begin{equation}
 h_0^{\mathrm{wt}}(\mathcal M,\mathcal Y^d)=0.
 \label{null:current-hom}
\end{equation}

Suppose $\mathcal Y^d$ has a socle in color $j\in J$.  Then the top
of $\mathcal M$ at that color vanishes.  Otherwise the proportionality
of tops across the full label family supplies a simple coefficient-array
quotient at the same unfolded vertex as the chosen socle, in some member
of $\mathcal M$.  Composing that quotient with the socle inclusion
gives a nonzero homogeneous length-zero cocycle.  Explicitly, a functional
killing the incoming coefficient images followed by a vector killed by
the outgoing coefficients has both differential compositions zero.
This contradicts
Equation~\eqref{null:current-hom}.  The top of $\mathcal M$ there is
the top of the full level-$q$ layer.  Thus $F_j$ preserves the cap by
Lemma~\ref{cap:access}.

\emph{Preservation of earlier stages.}
For this single move, compare the old and new data by ascending through
$k<d$.  Suppose the old top at $j$ in level $k-1$ has already been
shown to vanish; this is vacuous below the lower bound.  The adjacent
layer sequence in Lemma~\ref{lay:adjacent} identifies
\begin{equation}
 \mathcal Y^k_{\mathrm{new}}
    \cong Y^k\bigl(F_j\mathcal Y^k_{\mathrm{old}}\bigr).
 \label{null:unshifted}
\end{equation}
There is no contribution from the old level $k-1$, precisely because
its top at $j$ vanishes.  The zero socle on $J$ is preserved by this
unshifted operation, by Lemma~\ref{lay:positive}.

The new representative is capped and satisfies
Equation~\eqref{null:interior}.  Lemma~\ref{null:separation} therefore
gives all the new conditions in Equation~\eqref{null:induction} at
level $k$.  In particular the old and new Cartan rows at $j$ both
vanish.  Multiplicities outside $j$ are unchanged, and $C_{jj}=2$,
so $m^k_{\mathrm{new},j}=m^k_{\mathrm{old},j}$.  If $t_j^k\ge0$
is the folded weighted old top at $j$ in level $k$, the pure positive
layer formula gives
\[
 m^k_{\mathrm{new},j}
   =m^k_{\mathrm{old},j}-(Cm^k_{\mathrm{old}})_j+t_j^k.
\]
Hence $t_j^k=0$.  Positivity of all label weights makes every old
top in that color zero.  This supplies the premise for the next
level; in particular, it proves Equation~\eqref{null:unshifted} also
at $k=d$.  It also restores the preceding-stage conditions needed
to apply Lemma~\ref{null:injection} after this move.

This order avoids circularity: top vanishing at $k-1$ first identifies
the operation at $k$; separation and the row comparison then establish
top vanishing at $k$.

\emph{Termination and contradiction.}
Continue firing a color of $J$ whenever a socle remains at level $d$.
At that level the operations are the unshifted pure-layer operations.
By Lemma~\ref{lay:bottom}, they satisfy the positive braid relations,
and by Lemma~\ref{lay:positive} a last acted color has zero socle
after its operation.  Write the most recent letter on the left.
If the next chosen color were a left descent of the reduced Coxeter
word already fired, Matsumoto's Theorem would move that letter to
the left end, presenting the current layer as the result of its
operation acting last.  Its socle would then be zero, a contradiction.
Every choice is therefore a Coxeter ascent.  Sphericity of $J$ bounds
the number of choices by the length of its longest element.

When the firing stops, the socle on $J$ at level $d$ vanishes.
Lemma~\ref{null:separation} supplies Equation~\eqref{null:induction}
for this level and, if $d<d_b$, permits the next stage.  There are
only finitely many stages.  At $d=d_b$ the same lemma gives
$m^{d_b}|_{\partial I}=0$, contradicting its fixed occupied boundary
level.  If $J=\varnothing$, there are no firing steps: the hull has
its zero-dimensional meaning and the separation lemma already gives
the same contradiction.  This proves the proposition.
\end{proof}
